% !TeX root = ../main.tex

\chapter{Data, Cohorts and Preprocessing}\label{ch:data}

\section{The DELCODE Cohort and the Conversion Label}\label{sec:delcode}
\outlinenote{Cohort composition, and the label definition together with the follow-up
window that defines it. States the tradeoff explicitly: the window choice trades cohort
size against label noise.}

\section{Longitudinal Sampling and Follow-Up Regularity}\label{sec:sampling}
\outlinenote{Visit spacing, missing follow-ups, variable sequence length, and the masking
and time-conditioning consequences. Records that DELCODE follow-up is more regular than the
irregular-sampling literature assumes, and that the pooled protocol reintroduces
irregularity.}

\section{External Cohorts: ADNI and OASIS-3}\label{sec:external-cohorts}

As part of the broader project, \ac{ADNI} and \ac{OASIS} were included as additional
cohorts for cross-cohort validation. During the thesis, these datasets were processed and
integrated into the existing analysis framework. This section describes the acquisition
heterogeneity of both cohorts relative to \ac{DELCODE}, the diagnostic-label conventions
used to identify conversion in each, and a scan-to-visit matching problem specific to
multi-site data that does not arise in \ac{DELCODE}.

\ac{DELCODE} is acquired at a single site under one nominal protocol
(\autoref{sec:preprocessing} gives the exact scanning parameters), so within-cohort
acquisition variability is not a concern for it. \ac{ADNI} and \ac{OASIS}, by contrast,
pool scans collected over more than a decade across many sites, scanner vendors and
hardware generations, which is the acquisition-heterogeneity challenge introduced in
\autoref{ch:foundations}. An internal inventory of the \ac{ADNI} magnetic-resonance record
metadata used for this project counts 18,616 records, of which 18,614 are acquired at
3.0\,T field strength and only 2 at 1.5\,T, drawn from at least ten scanner models across
three vendors (Siemens, 14,056 records; Philips, 2,416; General Electric, 2,144). \ac{OASIS}
draws on a single vendor, Siemens, but spans six scanner generations over its acquisition
window; its record inventory totals 30,339 magnetic-resonance sessions, of which 5,114
scan rows across 2,138 sessions and 1,198 subjects are resting-state functional scans. Of
these, 205 rows are reference or calibration acquisitions (\texttt{SeriesDescription} of
\texttt{rsfmri\_ref}, repetition time 5.0\,s) rather than usable resting runs and are
excluded before any downstream matching.

\paragraph{Repetition-time grouping and a unit trap.} Cohort-level protocol comparisons
group scans by repetition time (below 1{,}000\,ms as one group, below 2{,}000\,ms as a
second, and 2{,}000\,ms and above as a third). \ac{OASIS} stores repetition and echo time
in seconds in its scan metadata (for example 2.2\,s, 0.027\,s), whereas \ac{ADNI} stores
the same fields directly in milliseconds; applying the grouping thresholds to \ac{OASIS}
without first multiplying by 1{,}000 places every scan in the wrong group. Once converted,
\ac{OASIS} resting-state scans populate only two of the three repetition-time groups: one
at 2{,}200, 2{,}500 or 3{,}000\,ms and a second at 500, 806 or 816\,ms; the middle group is
empty. This unit conversion was implemented as part of the cohort-integration work carried
out during the thesis, alongside the rest of the \ac{OASIS} subject-list construction
described below.

\paragraph{An era boundary in the diagnostic label.} \ac{ADNI} and \ac{OASIS} each label a
subject a converter from the same rule: a subject with at least one visit diagnosed
\ac{MCI} is a converter if any later visit is diagnosed \ac{AD}, and stable \ac{MCI}
otherwise; a subject whose only \ac{AD} diagnosis occurs at or before the first \ac{MCI}
visit is treated as prevalent \ac{AD} rather than a converter and is dropped, since the
label is meant to be predicted ahead of conversion rather than recognised in visits that
are already demented.
Applying this rule to \ac{OASIS} exposes a form-schema change that a same-cohort analysis
would not: \ac{OASIS} diagnoses are drawn from the Uniform Data Set clinical form, whose
Alzheimer's-disease field was renamed partway through the study. Earlier visits carry a
single field, \texttt{PROBAD}; later visits instead carry two fields, \texttt{alzdis} and
\texttt{alzdisif}, and leave \texttt{PROBAD} unset. Of 8{,}499 diagnostic rows, 4{,}855
belong to the earlier form era and 3{,}644 to the later one, and the two eras are
essentially non-overlapping: \texttt{PROBAD} is populated in 1{,}131 of the 4{,}855 earlier
rows and in none of the 3{,}644 later rows. A subject is therefore labelled \ac{AD} at a
visit if
\[
  \text{DEMENTED} = 1 \;\wedge\; \bigl(\text{PROBAD} = 1 \;\vee\; (\text{alzdis} = 1 \wedge \text{alzdisif} = 1)\bigr),
\]
where $\text{DEMENTED}$ is the form's general dementia flag, $\text{PROBAD}$ is the
earlier-era probable-\ac{AD} flag, and $\text{alzdis}$ together with $\text{alzdisif}$ are
the later-era Alzheimer's-disease and Alzheimer's-as-primary-cause flags that replace it.
Using $\text{PROBAD}$ alone, ignoring the later form era, yields 23 converters and 130
stable \ac{MCI} subjects; the era-aware rule above yields 73 converters and 85 stable
\ac{MCI} subjects, nearly tripling the converter count. The single-field rule under-counts
because it silently treats every later-era \ac{AD} visit as non-\ac{AD}, which is exactly
the failure mode that motivates checking a label rule against both form eras of a cohort
before trusting its counts, rather than assuming a diagnostic field keeps the same meaning
across a multi-year, multi-site study.

\paragraph{Matching a scan to a clinical visit.} Neither cohort scans a subject on the same
day as its clinical assessment, so a scan is assigned to the nearest visit within a fixed
window and left unmatched otherwise. \ac{ADNI} uses a window of $\pm 90$ days; \ac{OASIS}
uses $\pm 365$ days. The difference follows from visit cadence rather than from a different
tolerance for label noise: \ac{OASIS} clinical visits recur at approximately annual
intervals, so a $\pm 90$-day window matches only 33 visit-scan pairs across the cohort,
against 122 matched visit-scans over 84 subjects at $\pm 365$ days. A fixed matching window
carried over unchanged from \ac{ADNI} to \ac{OASIS} would therefore have discarded most of
the cohort's imaging data before any modelling step saw it.

\paragraph{No calendar dates in \ac{OASIS}.} \ac{ADNI} visit and scan records carry an
actual calendar exam date. \ac{OASIS}, being fully de-identified, carries no calendar date
anywhere in its metadata: every visit and scan is instead timestamped as an integer count
of days since an arbitrary per-subject entry point. The two cohorts' visit and scan
timestamp columns are given identical names for downstream compatibility, but the
\ac{OASIS} columns hold day offsets rather than dates, and parsing them as dates produces
meaningless values. This is a cohort-integration detail specific to processing \ac{OASIS},
addressed as part of the thesis when the \ac{OASIS} subject-list construction was built to
mirror the pre-existing \ac{ADNI} pipeline.

Applying the classification rule above to \ac{OASIS} yields 73 converters and 85 stable
\ac{MCI} subjects; 67 of the converters and 73 of the stable \ac{MCI} subjects have at
least one scan matched within the $\pm 365$-day window, for 122 matched visit-scans across
84 subjects in total. The corresponding \ac{ADNI} classification, built on the same rule
and applied to a substantially larger record inventory, is carried through the further
attrition and split-construction stages described in \autoref{sec:pooled-protocol} and
\autoref{sec:preprocessing}.

\section{The Pooled ADNI and DELCODE Protocol}\label{sec:pooled-protocol}
\outlinenote{The technical description of the asset-building pipeline.}

The pooled protocol implemented during the thesis builds on the pre-existing
\ac{DELCODE} split definitions and leakage-control framework, extending these structures
to the additional cohort.

\section{Preprocessing and Confound Management}\label{sec:preprocessing}
\outlinenote{Confound regression, motion handling and the attrition accounting.}

\ac{DELCODE} had been preprocessed before the start of this thesis using the laboratory's
existing preprocessing pipeline. During the thesis, the same established preprocessing
framework was applied to \ac{ADNI} and \ac{OASIS}. Cohort-specific adaptations required for
processing and integrating these datasets were implemented as part of the thesis.

\section{Group-Level Characterisation Before Modelling}\label{sec:group-stats}
\outlinenote{Converters against non-converters characterised with \ac{PERMANOVA} on the
connectivity matrices, the \ac{NBS} for cluster-level differences, and per-edge effect
sizes under \ac{FDR} control. Establishes whether a group difference is present before a
deep model is asked to find one.}

\section{Split Integrity and Leakage Control}\label{sec:leakage}
\outlinenote{Patient-level splitting, the isolated test manifest, post-conversion visit
exclusion, validation-derived thresholds, and the executable audit that enforces all of it.
The audit output is reproduced in \autoref{ch:app-splits}.}
